% Official solution (2012_da_sol.pdf, "Data analysis", August 10, 2012, Question 1).
% The published solution consists of a point distribution ("Pontuation", omitted here as a
% standalone marking scheme) and the answer list below; all quantitative results are
% presented in Figures 1-7.

Answers:
\begin{itemize}
\item Items 1,2 \textbf{See Fig. 1--3.}
\item Item 3: \textbf{Night B}
\item Items 4,5,6 \textbf{See Fig. 4--6.}
\item Items 7,8 \textbf{See Fig. 7.}
\item Item 9 \textbf{Yes}
\end{itemize}

% FIGURE NEEDED: Figure 1 (2012_da_sol.pdf p.1): instrumental magnitude vs air mass for
% Night A with fitted straight line, titled "Night A (A=0.3493)" -- extinction coefficient
% 0.3493.
% FIGURE NEEDED: Figure 2 (2012_da_sol.pdf p.1): same for Night B, titled
% "Night B (A=0.4285)".
% FIGURE NEEDED: Figure 3 (2012_da_sol.pdf p.2): same for Night C, titled
% "Night C (A=0.2059)".
% FIGURE NEEDED: Figure 4 (2012_da_sol.pdf p.2): calibrated magnitude vs time (hours) for
% Night A, with annotations "maxima interval = 2.5 h", "minima interval = 3.8 h",
% "Amplitude = 0.09".
% FIGURE NEEDED: Figure 5 (2012_da_sol.pdf p.2): calibrated magnitude vs time for Night B,
% with annotations "maxima interval = 3.1 h", "minima interval = 3.1 h", "Amplitude = 0.2".
% FIGURE NEEDED: Figure 6 (2012_da_sol.pdf p.3): calibrated magnitude vs time for Night C,
% with annotations "maxima interval = 3.3 h", "minima interval = 3.1 h", "Amplitude = 0.1".
% FIGURE NEEDED: Figure 7 (2012_da_sol.pdf p.3): reduced magnitude vs phase angle
% (degrees), titled "Phase curve (linear coef.: 0.0405)", straight-line fit through the
% three nights' points with the near-opposition point lying above the line.
